% ==========================================================================
%  Chapter 5 — UML Implementation Details
% ==========================================================================
\chapter{UML Implementation Details}
\label{ch:implementation}

\epigraph{%
  Show me the data structures and I won't usually need your code;
  it will be obvious.%
}{Fred Brooks, paraphrased}

This chapter descends from the architectural overview of the previous chapters
into the concrete data structures, code paths, and protocols that make
User-Mode Linux work.  Where earlier chapters described \emph{what} UML does,
this chapter describes \emph{how}: the exact sequence of system calls, signal
deliveries, futex operations, and page-table walks that execute when a guest
process issues a system call, takes a page fault, or is context-switched.
All code references are relative to the \texttt{arch/um/} subtree of the Linux
6.16+ source unless noted otherwise.

% ------------------------------------------------------------------
\section{Process Lifecycle}
\label{sec:impl-process-lifecycle}

A UML guest process is, from the host kernel's perspective, just another
userspace thread.  Its lifecycle maps onto the standard Linux process model
but with a UML-specific twist at each stage: creation, scheduling, and
destruction are all mediated through the active backend.

\subsection{Creation: \texttt{fork} through \texttt{copy\_thread}}
\label{sec:impl-process-creation}

When a guest process calls \syscall{fork} or \syscall{clone}, the standard
kernel path reaches \texttt{copy\_process()}, which eventually calls the
architecture-specific \texttt{copy\_thread()} in
\umlpath{arch/um/kernel/process.c}.  This function is the first point where
UML diverges from a hardware architecture port:

\begin{lstlisting}[language=KernelC,caption={Simplified \texttt{copy\_thread()} --- task creation dispatch},label=lst:copy-thread]
int copy_thread(struct task_struct *p,
                const struct kernel_clone_args *args)
{
    void (*handler)(void);

    p->thread = (struct thread_struct) INIT_THREAD;

    if (!args->fn) {
        /* Forking a user process: copy parent's register state */
        memcpy(&p->thread.regs.regs, current_pt_regs(),
               sizeof(p->thread.regs.regs));
        PT_REGS_SET_SYSCALL_RETURN(&p->thread.regs, 0);
        if (sp != 0)
            REGS_SP(p->thread.regs.regs.gp) = sp;
        handler = fork_handler;
        arch_copy_thread(&current->thread.arch, &p->thread.arch);
    } else {
        /* Kernel thread: set up fresh register state */
        um_backend_dispatch(init_thread_regs,
                            p->thread.regs.regs.gp,
                            p->thread.regs.regs.fp);
        p->thread.request.thread.proc = args->fn;
        p->thread.request.thread.arg  = args->fn_arg;
        handler = new_thread_handler;
    }

    um_backend_dispatch(thread_create, p,
                        task_stack_page(p), handler);
    return 0;
}
\end{lstlisting}

The key dispatch is \texttt{um\_backend\_dispatch(thread\_create, ...)} at the
end.  For the seccomp backend, this initializes the child's
\texttt{jmp\_buf} (the \texttt{switch\_buf} in \texttt{thread\_struct}) so
that the first context switch into the new task will \texttt{longjmp} to the
chosen \texttt{handler}---either \texttt{fork\_handler} for user processes or
\texttt{new\_thread\_handler} for kernel threads.

\subsection{The \texttt{thread\_struct} and Per-Task State}
\label{sec:impl-thread-struct}

Every Linux task has an architecture-specific \texttt{thread\_struct} embedded
in its \texttt{task\_struct}.  UML's definition
(\umlpath{arch/um/include/asm/processor-generic.h}) is compact but
information-dense:

\begin{lstlisting}[language=KernelC,caption={UML's \texttt{thread\_struct}},label=lst:thread-struct]
struct thread_struct {
    struct pt_regs     *segv_regs;     /* saved regs at SEGV time */
    struct task_struct *prev_sched;    /* previous task (for tail) */
    struct arch_thread  arch;          /* x86-specific: TLS, debug */
    jmp_buf             switch_buf;    /* longjmp target for ctxsw */
    struct {
        struct {
            int   (*proc)(void *);     /* kernel thread entry fn */
            void   *arg;               /* kernel thread argument */
        } thread;
    } request;
    void               *segv_continue;
    struct pt_regs      regs;          /* variable-sized: includes FP */
};
\end{lstlisting}

The \texttt{switch\_buf} field is the heart of the seccomp backend's context
switch mechanism: it is a \texttt{jmp\_buf} that records the exact point in
host userspace where this task's execution was last suspended.  The
\texttt{arch} sub-structure (\texttt{struct arch\_thread}) carries
architecture-dependent state such as TLS base addresses
(\texttt{fs\_base}/\texttt{gs\_base} on x86-64) and debug register state.

\subsection{Context Switch: \texttt{\_\_switch\_to}}
\label{sec:impl-context-switch}

The core scheduler's \texttt{context\_switch()} calls the architecture-specific
\texttt{\_\_switch\_to()}, which on UML dispatches to the backend:

\begin{lstlisting}[language=KernelC,caption={UML context switch},label=lst:switch-to]
struct task_struct *__switch_to(struct task_struct *from,
                               struct task_struct *to)
{
    to->thread.prev_sched = from;
    set_current(to);
    um_on_context_switch(from, to);
    um_backend_dispatch(context_switch, from, to);
    arch_switch_to(current);
    return current->thread.prev_sched;
}
\end{lstlisting}

For the \textbf{seccomp backend}, \texttt{context\_switch} performs a
\texttt{setjmp}/\texttt{longjmp}: the outgoing task saves its execution state
into \texttt{from->thread.switch\_buf} via \texttt{setjmp}, and then
\texttt{longjmp}s to the saved state in \texttt{to->thread.switch\_buf}.  This
is a purely userspace operation---no host kernel involvement, no system calls.
The host kernel never knows that UML has switched from one ``guest task'' to
another; the host thread simply jumps to a different code address.

For a future \textbf{KVM backend}, the context switch would instead
load the new task's register state into the VMCS/VMCB before the next
\texttt{KVM\_RUN} ioctl---the vCPU hardware itself enforces the switch.

\subsection{Destruction}

Task destruction follows the standard Linux \texttt{do\_exit()} path.  The
backend-specific cleanup is minimal for seccomp: the stub child process
(which is per-mm, not per-task) is destroyed only when the last task sharing
that mm exits, at which point \texttt{mm\_destroy()} is called through the
backend ops table.

% ------------------------------------------------------------------
\section{The Seccomp Backend: Trap Path Walkthrough}
\label{sec:impl-seccomp-trap}

The seccomp backend, implemented primarily in
\umlpath{arch/um/kernel/skas/stub.c} and
\umlpath{arch/um/os-Linux/skas/process.c}, is the production backend
for current UML.  This section walks through the complete lifecycle of a
guest system call, step by step.

\subsection{Step 1: Guest Issues a System Call}

A guest userspace process (running as a host thread inside a seccomp-filtered
address space) executes \texttt{syscall} (x86-64) or \texttt{int 0x80}.  The
host kernel's syscall entry path is reached, but \emph{before} the syscall
can execute, the BPF seccomp filter installed by UML's stub setup intercepts
it.

The filter returns \texttt{SECCOMP\_RET\_TRAP}, which causes the host kernel
to:
\begin{enumerate}[nosep]
  \item Abort the system call (it never executes on the host).
  \item Deliver a \texttt{SIGSYS} signal to the stub child process.
\end{enumerate}

\subsection{Step 2: \texttt{stub\_signal\_interrupt} --- The Stub Handler}
\label{sec:impl-stub-signal-interrupt}

The \texttt{SIGSYS} signal is caught by
\texttt{stub\_signal\_interrupt()}\linebreak (\umlpath{arch/um/kernel/skas/stub.c}),
which runs on the signal stack inside the stub child's address space.
This handler is the most performance-critical code in the seccomp backend;
it executes for every guest system call, signal, and timer interrupt.

The handler performs five critical operations:

\begin{enumerate}
\item \textbf{Save signal context.}  The signal number, \texttt{siginfo\_t}
  offset, and \texttt{mcontext\_t} offset (both relative to the signal stack
  base) are stored into the shared \texttt{stub\_data} page:
\begin{lstlisting}[language=KernelC,numbers=none]
d->signal    = sig;
d->si_offset = (unsigned long)info
             - (unsigned long)&d->sigstack[0];
d->mctx_offset = (unsigned long)&uc->uc_mcontext
               - (unsigned long)&d->sigstack[0];
\end{lstlisting}

\item \textbf{TLS base synchronization} (x86-64 only).  Because UML's
  seccomp filter \texttt{ALLOW}s \texttt{arch\_prctl}, guest userspace can
  change \texttt{FS\_BASE}/\texttt{GS\_BASE} directly without a SIGSYS
  round-trip.  The stub must re-read the authoritative values from the host
  kernel before handing control back to UML, or subsequent set/get comparisons
  will operate on stale data---a bug that manifested as Python segfaults in
  \texttt{\_Py\_Dealloc} through NULL TLS dereferences:
\begin{lstlisting}[language=KernelC,numbers=none]
#ifdef CONFIG_X86_64
stub_syscall2(__NR_arch_prctl, ARCH_GET_FS,
              (unsigned long)&d->arch_data.fs_base);
stub_syscall2(__NR_arch_prctl, ARCH_GET_GS,
              (unsigned long)&d->arch_data.gs_base);
#endif
\end{lstlisting}

\item \textbf{Wake the UML kernel via futex.}  The stub sets
  \texttt{d->futex = FUTEX\_IN\_KERN} and issues a
  \texttt{FUTEX\_WAKE} on the futex word, then enters a
  \texttt{FUTEX\_WAIT} loop until the UML kernel changes the futex value
  back (see Section~\ref{sec:impl-futex-protocol}).

\item \textbf{Execute queued stub syscalls} (if any).  When the UML kernel
  needs to modify the stub child's address space (mmap, munmap), it queues
  the operations in \texttt{stub\_data.syscall\_data[]} and wakes the stub.
  The stub processes these via the \texttt{syscall\_handler()} inline function,
  receiving any file descriptors via SCM\_RIGHTS over a Unix socket.

\item \textbf{Restore architecture-dependent state.}  Before returning from
  the signal handler (which triggers \texttt{rt\_sigreturn} to restore the
  guest's mcontext), the stub calls
  \texttt{stub\_seccomp\_restore\_state(\&d->arch\_data)} to apply any
  register modifications the UML kernel wrote into the shared page.
\end{enumerate}

\begin{figure}[ht]
\centering
\begin{tikzpicture}[
    box/.style={
      draw=UMLBlue!70, fill=UMLBlue!5, rounded corners=3pt,
      minimum width=3.6cm, minimum height=0.9cm,
      font=\small, align=center, text=UMLInk
    },
    stubbox/.style={
      draw=UMLTeal!70, fill=UMLTeal!5, rounded corners=3pt,
      minimum width=3.6cm, minimum height=0.9cm,
      font=\small, align=center, text=UMLInk
    },
    hostbox/.style={
      draw=UMLMuted!70, fill=UMLPanel, rounded corners=3pt,
      minimum width=3.6cm, minimum height=0.9cm,
      font=\small, align=center, text=UMLInk
    },
    arr/.style={-{Stealth[length=5pt]}, thick, color=UMLBlue!70},
    darr/.style={-{Stealth[length=5pt]}, thick, color=UMLTeal!70, dashed},
    note/.style={font=\scriptsize\itshape, text=UMLMuted},
    every node/.style={inner sep=4pt}
  ]

  % Guest side
  \node[box] (guest) at (0, 0) {Guest calls\\[-2pt]\texttt{write(2)}};

  % Host kernel
  \node[hostbox] (seccomp) at (0, -1.6) {Host kernel:\\[-2pt]seccomp filter\\[-2pt]\texttt{RET\_TRAP}};

  % Stub handler
  \node[stubbox] (sigsys) at (0, -3.2) {Stub receives\\[-2pt]\texttt{SIGSYS}};
  \node[stubbox] (save) at (0, -4.8) {Save regs to\\[-2pt]\texttt{stub\_data}};
  \node[stubbox] (tls) at (0, -6.4) {Sync TLS:\\[-2pt]\texttt{ARCH\_GET\_FS/GS}};
  \node[stubbox] (futex1) at (0, -8.0) {\texttt{FUTEX\_WAKE}\\[-2pt]UML kernel};

  % UML kernel side
  \node[box] (dispatch) at (5.8, -8.0) {UML kernel:\\[-2pt]\texttt{handle\_syscall}};
  \node[box] (result) at (5.8, -6.4) {Write result to\\[-2pt]\texttt{stub\_data}};
  \node[box] (futex2) at (5.8, -4.8) {\texttt{FUTEX\_WAKE}\\[-2pt]stub child};

  % Return path
  \node[stubbox] (restore) at (0, -9.6) {Restore state;\\[-2pt]\texttt{rt\_sigreturn}};
  \node[box] (resume) at (0, -11.2) {Guest resumes\\[-2pt]at next insn};

  % Arrows
  \draw[arr] (guest) -- (seccomp);
  \draw[arr] (seccomp) -- (sigsys);
  \draw[arr] (sigsys) -- (save);
  \draw[arr] (save) -- (tls);
  \draw[arr] (tls) -- (futex1);
  \draw[darr] (futex1) -- (dispatch) node[note, midway, above] {futex wake};
  \draw[arr] (dispatch) -- (result);
  \draw[arr] (result) -- (futex2);
  \draw[darr] (futex2) -- (restore) node[note, midway, left, xshift=-3mm] {futex wake};
  \draw[arr] (restore) -- (resume);

  % Annotations
  \node[note, anchor=east] at (-2.2, -4.8) {stub address space};
  \node[note, anchor=west] at (8.0, -6.4) {UML kernel process};

  % Bracket for stub side
  \draw[decorate, decoration={brace, amplitude=6pt, mirror},
        UMLTeal!50, thick]
    (-2.0, -2.8) -- (-2.0, -10.0)
    node[note, midway, left, xshift=-8pt, align=right]
    {Runs inside\\stub child};

\end{tikzpicture}
\caption{The seccomp trap path: complete lifecycle of a guest system call.
Solid arrows show the control flow; dashed arrows show the futex-mediated
handoff between the stub child and the UML kernel process.}
\label{fig:seccomp-trap-path}
\end{figure}

\subsection{The Futex Protocol}
\label{sec:impl-futex-protocol}

Communication between the stub child and the UML kernel uses a shared-memory
futex protocol~\cite{futex-lwn}.  The futex word lives at
\texttt{stub\_data.futex} and takes two values:

\begin{center}
\begin{tabular}{lll}
\toprule
\textbf{Constant} & \textbf{Value} & \textbf{Meaning} \\
\midrule
\texttt{FUTEX\_IN\_CHILD} & 0 & Stub child is running (guest mode) \\
\texttt{FUTEX\_IN\_KERN}  & 1 & UML kernel owns control \\
\bottomrule
\end{tabular}
\end{center}

The protocol proceeds as follows:

\begin{enumerate}[nosep]
\item \textbf{Stub $\to$ kernel:} The stub sets
  \texttt{futex = FUTEX\_IN\_KERN}, then issues
  \texttt{futex(FUTEX\_WAKE, 1)} to wake the UML kernel thread.
\item \textbf{Stub waits:} The stub loops on
  \texttt{futex(FUTEX\_WAIT, FUTEX\_IN\_KERN)} until the value changes.
  Spurious wakeups (\texttt{-EINTR}) are retried; unexpected errors cause
  the stub to \texttt{exit\_group(1)}.
\item \textbf{Kernel processes:} The UML kernel reads the saved registers and
  signal information from \texttt{stub\_data}, dispatches to
  \texttt{handle\_syscall()}, writes the result back, and then sets
  \texttt{futex = FUTEX\_IN\_CHILD} followed by
  \texttt{futex(FUTEX\_WAKE, 1)}.
\item \textbf{Stub resumes:} The stub's \texttt{FUTEX\_WAIT} returns; it
  checks for queued syscalls and, finding none (or having processed them),
  calls \texttt{stub\_seccomp\_restore\_state()} and returns from the signal
  handler, which triggers \texttt{rt\_sigreturn} and resumes guest execution.
\end{enumerate}

The critical invariant is that only one side is ``active'' at any time: when
the futex word is \texttt{FUTEX\_IN\_KERN}, the stub is sleeping in the
kernel's futex wait queue and will not execute any guest code; when it is
\texttt{FUTEX\_IN\_CHILD}, the UML kernel thread is sleeping.  This mutual
exclusion is what allows the shared \texttt{stub\_data} page to be accessed
without locks.

\subsection{The \texttt{stub\_data} Structure}
\label{sec:impl-stub-data}

The shared communication page between the stub child and the UML kernel is
defined in \umlpath{arch/um/include/shared/skas/stub-data.h}:

\begin{lstlisting}[language=KernelC,caption={The \texttt{stub\_data} structure},label=lst:stub-data-impl]
struct stub_data {
    long err;                    /* syscall error from stub handler */

    int syscall_data_len;        /* number of queued stub syscalls */
    /* Batched syscall descriptors, sized to fill one page minus
     * overhead for the other fields (128 bytes reserved) */
    struct stub_syscall syscall_data[...]  __aligned(16);

    /* Shared with signal handler (seccomp mode only) */
    short         restart_wait;  /* re-enter wait loop after batch */
    unsigned int  futex;         /* FUTEX_IN_CHILD / FUTEX_IN_KERN */
    int           signal;        /* signal number (e.g., SIGSYS) */
    unsigned short si_offset;    /* siginfo_t offset in sigstack */
    unsigned short mctx_offset;  /* mcontext_t offset in sigstack */

    /* Architecture-specific state for seccomp restore */
    struct stub_data_arch arch_data;

    /* Signal stack for the stub's signal handlers */
    unsigned char sigstack[UM_KERN_PAGE_SIZE]
                  __aligned(UM_KERN_PAGE_SIZE);
};
\end{lstlisting}

The \texttt{syscall\_data[]} array carries batched memory operations:

\begin{lstlisting}[language=KernelC,caption={Stub syscall descriptor},label=lst:stub-syscall]
enum stub_syscall_type {
    STUB_SYSCALL_UNSET = 0,
    STUB_SYSCALL_MMAP,
    STUB_SYSCALL_MUNMAP,
};

struct stub_syscall {
    struct {
        unsigned long addr;
        unsigned long length;
        unsigned long offset;
        int fd;
        int prot;
    } mem;
    enum stub_syscall_type syscall;
};
\end{lstlisting}

The batching mechanism is crucial for TLB synchronization performance: instead
of waking the stub for each individual mmap/munmap, the UML kernel can queue
multiple operations and wake the stub once.

% ------------------------------------------------------------------
\section{Memory Management Implementation}
\label{sec:impl-memory}

UML's memory subsystem must bridge two page-table hierarchies: the guest page
tables that the Linux mm subsystem manages, and the host page tables that the
host kernel manages for the stub child's address space.  The bridge between
them is the TLB synchronization path.

\subsection{Physical Memory: The Backing File}
\label{sec:impl-physmem}

All guest physical memory is backed by a single temporary file, created
at boot by \texttt{setup\_physmem()} in \umlpath{arch/um/kernel/physmem.c}.
The function calls \texttt{create\_mem\_file(len)} to create an unlinked
temporary file of size \texttt{physmem\_size}, then \texttt{mmap}s it into
the UML process's address space with \texttt{MAP\_SHARED}.  This file
descriptor---stored in the static \texttt{physmem\_fd}---is the single source
of truth for all guest RAM.

The \texttt{phys\_mapping()} function translates a guest physical address to
a (file descriptor, offset) pair:

\begin{lstlisting}[language=KernelC,caption={Physical address to file mapping},label=lst:phys-mapping]
int phys_mapping(unsigned long phys, unsigned long long *offset_out)
{
    int fd = -1;
    if (phys < physmem_size) {
        fd = physmem_fd;
        *offset_out = phys;
    }
    return fd;
}
\end{lstlisting}

When the seccomp backend needs to map a guest page into the stub child's
address space, it issues an \texttt{mmap(MAP\_SHARED | MAP\_FIXED)} on
\texttt{physmem\_fd} at the appropriate offset.  Because both the UML kernel
and the stub child map the same file with \texttt{MAP\_SHARED}, changes made
by the kernel (e.g., writing to a page that backs a guest buffer) are
immediately visible to the stub child, and vice versa.

\subsection{The TLB Sync Path: \texttt{um\_tlb\_sync}}
\label{sec:impl-tlb-sync}

When the guest kernel modifies a page table entry (e.g., through
\texttt{set\_pte}, \texttt{set\_ptes}, or any of the \texttt{flush\_tlb\_*}
family), UML does not immediately reflect the change in the host address space.
Instead, it \emph{marks} the affected range for deferred synchronization:

\begin{lstlisting}[language=KernelC,caption={TLB mark for deferred sync (from \texttt{pgtable.h})},label=lst:tlb-mark]
static inline void um_tlb_mark_sync(struct mm_struct *mm,
                                    unsigned long start,
                                    unsigned long end)
{
    /* Expand the pending sync range to include [start, end) */
    if (!mm->context.sync_tlb_range_to ||
         start < mm->context.sync_tlb_range_from)
        mm->context.sync_tlb_range_from = start;
    if (end > mm->context.sync_tlb_range_to)
        mm->context.sync_tlb_range_to = end;
}
\end{lstlisting}

The actual synchronization happens in \texttt{um\_tlb\_sync()}
(\umlpath{arch/um/kernel/tlb.c}), which is called before the guest is allowed
to resume execution---typically from the backend's \texttt{vcpu\_run} dispatch.
The function walks the guest page directory over the marked range
\texttt{[sync\_tlb\_range\_from, sync\_tlb\_range\_to)}, inspecting each leaf
PTE:

\begin{itemize}[nosep]
\item If the PTE is \textbf{present} and \texttt{pte\_needsync()} returns true,
  the function constructs a \texttt{struct um\_memory\_region} and dispatches
  through the backend's \texttt{mm\_region\_added} op (which, for seccomp,
  queues a \texttt{STUB\_SYSCALL\_MMAP} into \texttt{stub\_data}).
\item If the PTE is \textbf{not present} and \texttt{pte\_needsync()} returns
  true, a \texttt{mm\_region\_removed} dispatch queues a
  \texttt{STUB\_SYSCALL\_MUNMAP}.
\item On success, the PTE is marked up-to-date via \texttt{pte\_mkuptodate()}.
\end{itemize}

\subsection{The \texttt{mm\_context\_t} and Generation Tracking}
\label{sec:impl-mm-context}

UML extends the standard \texttt{mm\_struct} with a \texttt{mm\_context\_t}
(\umlpath{arch/um/include/asm/mmu.h}) that carries per-mm state for the
TLB sync and backend interaction:

\begin{lstlisting}[language=KernelC,caption={UML's \texttt{mm\_context\_t} (key fields)},label=lst:mm-context]
typedef struct mm_context {
    struct mm_id    id;              /* stub child PID, socket, etc. */
    struct mutex    turnstile;       /* serialize stub access */

    /* Deferred TLB sync range */
    spinlock_t      sync_tlb_lock;
    unsigned long   sync_tlb_range_from;
    unsigned long   sync_tlb_range_to;

    /* Deferred page free (mmu_gather integration) */
    spinlock_t      deferred_free_lock;
    struct list_head deferred_free_pages;
    unsigned int    deferred_free_count;

    /* Per-mm TLB generation counter */
    atomic64_t      tlb_gen;
    atomic64_t      tlb_gen_seen_by[NR_CPUS];
} mm_context_t;
\end{lstlisting}

The \texttt{mm\_id} sub-structure identifies the stub child associated with
this mm:

\begin{lstlisting}[language=KernelC,caption={The \texttt{mm\_id} structure},label=lst:mm-id]
struct mm_id {
    int pid;                     /* stub child PID */
    unsigned long stack;         /* shared stub_data page address */
    int syscall_data_len;        /* current batch length */
    /* Seccomp-only fields */
    int sock;                    /* Unix socket to stub child */
    int syscall_fd_num;          /* number of fds to send */
    int syscall_fd_map[STUB_MAX_FDS]; /* fd table for SCM_RIGHTS */
};
\end{lstlisting}

The \texttt{tlb\_gen} counter is incremented after each successful
\texttt{um\_tlb\_sync()} drain.  Under SMP configurations, each vCPU tracks
\texttt{tlb\_gen\_seen\_by[cpu]}; when a vCPU's seen generation lags behind the
current generation, it knows a TLB flush is needed before resuming guest
execution.

\subsection{The \texttt{um\_memory\_region} Descriptor}
\label{sec:impl-memory-region}

The interface between the mm-arbiter (TLB sync) and the backend uses a
clean, backend-agnostic descriptor:

\begin{lstlisting}[language=KernelC,caption={Memory region descriptor},label=lst:um-memory-region]
struct um_memory_region {
    unsigned long va;        /* guest virtual address */
    unsigned long len;       /* region length */
    int           prot;      /* UM_PROT_READ|WRITE|EXEC */
    int           phys_fd;   /* physmem fd (-1 for unmap) */
    u64           offset;    /* offset into phys_fd */
    void         *backend_data;  /* opaque, backend-managed */
};
\end{lstlisting}

For the seccomp backend, \texttt{mm\_region\_added} translates this into a
\texttt{mmap(MAP\_SHARED | MAP\_FIXED, phys\_fd, offset)} in the stub child's
address space.  The \texttt{backend\_data} field is reserved for the KVM
backend, which will use it to track memslot IDs.

% ------------------------------------------------------------------
\section{Signal Handling}
\label{sec:impl-signals}

UML repurposes the host's POSIX signal mechanism to implement kernel
``interrupts.''  Three signals carry the bulk of the work:

\begin{center}
\begin{tabular}{lll}
\toprule
\textbf{Host Signal} & \textbf{UML Semantics} & \textbf{Handler} \\
\midrule
\texttt{SIGALRM} & Timer interrupt (HZ tick) & \texttt{timer\_alarm\_handler} \\
\texttt{SIGIO}   & I/O event notification    & \texttt{sig\_handler} $\to$ \texttt{sigio\_handler} \\
\texttt{SIGSEGV} & Page fault                & \texttt{sig\_handler} $\to$ \texttt{segv\_handler} \\
\texttt{SIGCHLD} & Stub child death (seccomp) & \texttt{sig\_handler} $\to$ \texttt{sigchld\_handler} \\
\bottomrule
\end{tabular}
\end{center}

\subsection{Signal Installation and Masking}
\label{sec:impl-signal-masking}

All UML signal handlers are installed by \texttt{set\_handler()} in
\umlpath{arch/um/os-Linux/signal.c}, which uses \texttt{sigaction} with
\texttt{SA\_SIGINFO | SA\_ONSTACK}.  The function constructs a signal mask
that blocks all IRQ-class signals (\texttt{SIGIO}, \texttt{SIGWINCH},
\texttt{SIGALRM}) inside any handler, preventing nested interrupts.
\texttt{SIGCHLD} is conditionally masked only when the active backend
uses the stub reaper (\texttt{um\_backend->uses\_stub\_reaper}):

\begin{lstlisting}[language=KernelC,caption={Conditional SIGCHLD masking based on backend capability},label=lst:sigchld-mask]
if (um_backend && um_backend->uses_stub_reaper)
    sigaddset(&action.sa_mask, SIGCHLD);
\end{lstlisting}

\texttt{SIGSEGV} is installed with \texttt{SA\_NODEFER} so that nested
page faults can be handled (the page fault handler may itself access
unmapped memory during its processing).  Similarly, \texttt{SIGTRAP} uses
\texttt{SA\_NODEFER} when kprobes are enabled, since a probe handler can call
a function that is itself probed.

\subsection{Signals as Interrupts: The Enable/Disable Discipline}
\label{sec:impl-signal-discipline}

UML implements the kernel's \texttt{local\_irq\_disable()}/\texttt{local\_irq\_enable()} as signal blocking, not via
\texttt{sigprocmask} (which would require a system call on every
enable/disable), but through a per-thread software flag:

\begin{lstlisting}[language=KernelC,caption={Software-based IRQ disable},label=lst:signals-enabled]
__thread int signals_enabled;
static __thread unsigned int signals_pending;

void block_signals(void) {
    if (!signals_enabled) return;
    signals_enabled = 0;
    barrier();
}
\end{lstlisting}

When a signal arrives with \texttt{signals\_enabled == 0}, the handler records
the signal in \texttt{signals\_pending} and returns immediately without
processing it.  When \texttt{unblock\_signals()} is later called, it checks
\texttt{signals\_pending} and dispatches any accumulated signals in priority
order: \texttt{SIGIO} first (to avoid stranding I/O events), then
\texttt{SIGCHLD}, then \texttt{SIGALRM}:

\begin{lstlisting}[language=KernelC,caption={Deferred signal dispatch in \texttt{unblock\_signals()}},label=lst:unblock-dispatch]
void unblock_signals(void) {
    /* ... enable signals ... */
    while (1) {
        barrier();
        save_pending = signals_pending;
        if (save_pending == 0) return;
        signals_pending = 0;
        __block_signals();

        if (save_pending & SIGIO_MASK)
            sig_handler_common(SIGIO, NULL, NULL);
        if (save_pending & SIGCHLD_MASK)
            sigchld_handler(SIGCHLD, NULL, &regs, NULL);
        if ((save_pending & SIGALRM_MASK) &&
            (!(signals_active & SIGALRM_MASK)))
            timer_real_alarm_handler(NULL);

        __unblock_signals();
    }
}
\end{lstlisting}

This loop structure means that signals can ``stack'': a timer interrupt that
arrives while processing an I/O event will be caught on the next iteration.
The \texttt{signals\_active} bitmask prevents reentrant timer handling.

\subsection{The Timer Interrupt}

\texttt{SIGALRM} drives UML's tick.  The backend's \texttt{set\_timer} op
configures the host timer (via \texttt{setitimer} or \texttt{timer\_settime})
in one-shot or periodic mode.  When the alarm fires,
\texttt{timer\_alarm\_handler()} blocks signals, calls
\texttt{timer\_real\_alarm\_handler()} (which invokes the standard kernel
\texttt{timer\_handler}), and re-enables signals.  The standard kernel timer
subsystem then drives the scheduler tick, process accounting, and all
timer-wheel-based timeouts.

% ------------------------------------------------------------------
\section{The hostfs Filesystem}
\label{sec:impl-hostfs}

The \texttt{hostfs} filesystem (\umlpath{fs/hostfs/}) provides UML guests
with direct access to the host's filesystem tree.  It implements the Linux VFS
interface by forwarding each operation to the corresponding host system call.

\subsection{Architecture}

\texttt{hostfs} is structured as a two-layer filesystem:

\begin{itemize}[nosep]
\item \textbf{Kernel layer} (\texttt{hostfs\_kern.c}): Implements VFS
  operations (\texttt{inode\_operations}, \texttt{file\_operations},
  \texttt{address\_space\_operations}) using the internal API.
\item \textbf{User layer} (\texttt{hostfs\_user.c}): Implements the actual
  host system calls (\texttt{open}, \texttt{read}, \texttt{write},
  \texttt{stat}, etc.) in USER-mode compilation units that can include
  host \texttt{<fcntl.h>}, \texttt{<sys/stat.h>}, etc.
\end{itemize}

Each \texttt{hostfs\_inode\_info} carries an open host file descriptor
(\texttt{fd}) alongside the standard \texttt{vfs\_inode}.  Guest read/write
operations are forwarded directly to the host fd.  The mount option specifies
which host directory serves as the root of the mounted filesystem.

\subsection{io\_uring Integration for Async Writeback}

Recent UML adds an optional \texttt{io\_uring} substrate for batched
asynchronous writeback.  The per-superblock \texttt{hostfs\_fs\_info} carries
a \texttt{struct os\_io\_ring *wb\_ring} created at mount time with a depth of
128 entries.  When available, the writepages path submits write SQEs to the
ring instead of issuing synchronous \texttt{write} calls, improving throughput
for workloads with heavy file I/O.  If io\_uring is unavailable (older host
kernel, seccomp filtering blocks it), the path falls back transparently to
the legacy synchronous \texttt{write\_file()} loop.

\subsection{Security: openat2 and RESOLVE\_BENEATH}

When the \texttt{resolve\_strict} mount option is active, \texttt{hostfs}
routes file opens through \texttt{openat2()} with
\texttt{RESOLVE\_BENEATH | RESOLVE\_NO\_MAGICLINKS}, using an \texttt{O\_PATH}
directory fd opened on the mount root.  This prevents symlink-escape attacks
where a guest process creates a symlink pointing outside the configured
hostfs root.

% ------------------------------------------------------------------
\section{Boot Sequence}
\label{sec:impl-boot}

UML boots as a regular userspace executable.  The sequence from
\texttt{./linux} invocation to a running guest is shown in
Figure~\ref{fig:boot-sequence} and proceeds through these stages.

\subsection{Stage 1: Host Process Setup (\texttt{os-Linux/main.c})}
\label{sec:impl-boot-stage1}

The UML binary's \texttt{main()} performs host-level initialization:

\begin{enumerate}[nosep]
\item \textbf{Disable ASLR.} Calls \texttt{personality(ADDR\_NO\_RANDOMIZE)}
  and re-execs itself if the personality changed.  UML requires deterministic
  address layout for its stub mappings.
\item \textbf{Set stack limit.} Caps \texttt{RLIMIT\_STACK} at 8\,MB to
  ensure consistent stack placement.
\item \textbf{Install fatal handlers.} Registers \texttt{SIGINT} and
  \texttt{SIGTERM} handlers that call \texttt{\_exit(1)} immediately
  (async-signal-safe).
\item \textbf{Enter \texttt{linux\_main()}.}  Control transfers to
  \umlpath{arch/um/kernel/um\_arch.c}.
\end{enumerate}

\subsection{Stage 2: Kernel Pre-Boot (\texttt{linux\_main})}
\label{sec:impl-boot-stage2}

\texttt{linux\_main()} bridges host userspace and the kernel boot path:

\begin{enumerate}[nosep]
\item \textbf{Parse command line.}  Iterates \texttt{argv}, calling
  \texttt{uml\_checksetup()} for UML-specific parameters (\texttt{mem=},
  \texttt{root=}, \texttt{backend=}, etc.).
\item \textbf{Compute address space layout.}  Determines
  \texttt{host\_task\_size}, \texttt{stub\_start}, \texttt{task\_size},
  \texttt{physmem\_size}, VMALLOC boundaries.
\item \textbf{Run early host checks.}  \texttt{os\_early\_checks()} probes
  for seccomp, ptrace capabilities, and \texttt{/dev/kvm} availability.  Sets
  the \texttt{using\_seccomp} flag if the host supports seccomp-BPF.
\item \textbf{Select and initialize the backend.}  Calls
  \texttt{init\_backend()}, which selects the compiled-in backend (seccomp,
  KVM, or dynamically chosen), validates its HOT ops, runs \texttt{probe()}
  and \texttt{init()}, and logs the selection.
\item \textbf{Set up physical memory.}  Calls \texttt{setup\_physmem()},
  which creates the backing tempfile and maps it into the process.
\item \textbf{Enter \texttt{start\_uml()}.}  Calls
  \texttt{start\_kernel()}, the standard Linux boot entry point.
\end{enumerate}

\subsection{Stage 3: Standard Linux Boot (\texttt{start\_kernel})}
\label{sec:impl-boot-stage3}

From \texttt{start\_kernel()} onward, UML follows the same boot path as every
other Linux architecture port: \texttt{setup\_arch()} initializes
architecture-specific data structures, the mm subsystem initializes,
the scheduler starts, and eventually \texttt{kernel\_init()} mounts the root
filesystem and runs \texttt{/sbin/init}.

The root filesystem can be:
\begin{itemize}[nosep]
\item A \texttt{ubd} (User-Mode Block Device) backed by a host disk image.
\item A \texttt{hostfs} mount pointing at a host directory tree.
\item An \texttt{initrd} loaded from a host file.
\end{itemize}

\begin{figure}[ht]
\centering
\begin{tikzpicture}[
    stage/.style={
      draw=UMLBlue!70, fill=UMLBlue!5, rounded corners=3pt,
      minimum width=11cm, minimum height=1.0cm,
      font=\small, align=left, text=UMLInk,
      inner xsep=8pt
    },
    arr/.style={-{Stealth[length=5pt]}, thick, color=UMLBlue!50},
    note/.style={font=\scriptsize, text=UMLMuted, anchor=west}
  ]

  \node[stage] (s1) at (0, 0)
    {\textbf{1.}~\texttt{main()} ---
     Disable ASLR, set stack limit, install fatal handlers};
  \node[stage] (s2) at (0, -1.5)
    {\textbf{2.}~\texttt{linux\_main()} ---
     Parse cmdline, compute layout, \texttt{os\_early\_checks()}};
  \node[stage] (s3) at (0, -3.0)
    {\textbf{3.}~\texttt{init\_backend()} ---
     Select backend (seccomp/KVM), validate HOT ops, probe \& init};
  \node[stage] (s4) at (0, -4.5)
    {\textbf{4.}~\texttt{setup\_physmem()} ---
     Create tmpfile, mmap guest RAM (\texttt{MAP\_SHARED})};
  \node[stage] (s5) at (0, -6.0)
    {\textbf{5.}~\texttt{start\_kernel()} ---
     Standard Linux boot: mm init, scheduler, driver init};
  \node[stage] (s6) at (0, -7.5)
    {\textbf{6.}~\texttt{kernel\_init()} ---
     Mount root (ubd/hostfs), run \texttt{/sbin/init}};

  \draw[arr] (s1) -- (s2);
  \draw[arr] (s2) -- (s3);
  \draw[arr] (s3) -- (s4);
  \draw[arr] (s4) -- (s5);
  \draw[arr] (s5) -- (s6);

  \node[note] at (6.2, 0)     {\texttt{os-Linux/main.c}};
  \node[note] at (6.2, -1.5)  {\texttt{kernel/um\_arch.c}};
  \node[note] at (6.2, -3.0)  {\texttt{kernel/backend.c}};
  \node[note] at (6.2, -4.5)  {\texttt{kernel/physmem.c}};
  \node[note] at (6.2, -6.0)  {\texttt{init/main.c}};
  \node[note] at (6.2, -7.5)  {\texttt{init/main.c}};

\end{tikzpicture}
\caption{UML boot sequence from \texttt{./linux} invocation to running init.}
\label{fig:boot-sequence}
\end{figure}

% ------------------------------------------------------------------
\section{The Backend Ops Table}
\label{sec:impl-backend-ops}

The backend abstraction is the central architectural element of the
post-redesign UML.  All backend-specific behavior is accessed through a single
vtable: \texttt{struct um\_backend\_ops}, defined in
\umlpath{arch/um/include/shared/backend.h}.

\subsection{Structure Definition}

\begin{lstlisting}[language=KernelC,caption={The \texttt{um\_backend\_ops} structure (annotated excerpt)},label=lst:backend-ops-impl]
struct um_backend_ops {
    const char              *name;
    enum um_backend_kind     kind;
    u32                      contract_version;

    /* Capability flags */
    bool  uses_stub_reaper;        /* SIGCHLD for stub child death */
    bool  has_syscall_stub_fd_map; /* SCM_RIGHTS fd-passing */
    bool  stub_syscall_uses_futex; /* futex vs ptrace dispatch */
    bool  stub_child_runs_seccomp; /* SIGSYS filter in stub */

    /* Lifecycle and trap (4 ops) */
    int  (*probe)(void);
    int  (*init)(const struct um_backend_args *args);
    void (*shutdown)(void);
    void (*vcpu_run)(struct uml_pt_regs *regs);       /* HOT */

    /* Memory (5 ops) */
    int  (*mm_create)(struct mm_struct *mm);
    void (*mm_destroy)(struct mm_struct *mm);
    int  (*mm_region_added)(...);                      /* HOT */
    int  (*mm_region_removed)(...);                    /* HOT */
    int  (*mm_region_protected)(...);

    /* Scheduling (4 ops) */
    int  (*thread_create)(struct task_struct *p, ...);
    int  (*thread_start_idle)(void *stack, ...);
    void (*context_switch)(struct task_struct *prev,    /* HOT */
                           struct task_struct *next);
    int  (*ipi_send)(int cpu, int vector);

    /* Cross-vCPU TLB flush kick (may be NULL) */
    void (*tlb_kick_others)(struct mm_struct *mm);

    /* Time (3 ops) */
    u64  (*read_clock_ns)(void);                       /* HOT */
    int  (*set_timer)(int cpu, u64 deadline_ns,
                      enum um_timer_mode mode);
    u64  (*read_persistent_clock_ns)(void);

    /* Debug / introspection (3 ops) */
    void (*init_thread_regs)(unsigned long *gp,
                             unsigned long *fp);
    int  (*read_guest_regs)(struct task_struct *t, ...);
    int  (*write_guest_regs)(struct task_struct *t, ...);
};
\end{lstlisting}

\subsection{Dispatch Mechanism}

The dispatch macro in \texttt{backend.h} provides zero-cost abstraction for
single-backend builds and dynamic dispatch for multi-backend builds:

\begin{lstlisting}[language=KernelC,caption={Backend dispatch macro},label=lst:dispatch-macro-impl]
#if defined(CONFIG_UM_BACKEND_SECCOMP_ONLY)
  /* Direct call -- no indirect branch, no pointer dereference */
# define um_backend_dispatch(op, ...) \
         seccomp_##op(__VA_ARGS__)
#else
  /* Dynamic dispatch through the ops table pointer */
# define um_backend_dispatch(op, ...) \
         (um_backend->op(__VA_ARGS__))
#endif
\end{lstlisting}

In \texttt{SECCOMP\_ONLY} builds (the default configuration), every
\texttt{um\_backend\_dispatch(context\_switch, prev, next)} expands to
\texttt{seccomp\_context\_switch(prev, next)}---a direct function call with
no indirect branch prediction penalty.  In \texttt{DYNAMIC} builds, the
compiler generates an indirect call through the \texttt{um\_backend} pointer,
which is set once at boot by \texttt{init\_backend()} and never modified
afterward.

\subsection{Backend Selection at Boot}

\texttt{init\_backend()} in \umlpath{arch/um/kernel/backend.c} resolves the
backend at boot time through a three-step process:

\begin{enumerate}[nosep]
\item \textbf{Parse boot parameters.}  The \texttt{backend=} kernel command
  line parameter is parsed by \texttt{uml\_backend\_config()} into
  \texttt{backend\_arg\_requested} and \texttt{backend\_arg\_force}.
\item \textbf{Select backend.}  In \texttt{SECCOMP\_ONLY} builds, seccomp is
  always selected.  In \texttt{DYNAMIC} builds, \texttt{pick\_dynamic\_backend()}
  consults the boot parameter, falls back to seccomp if the requested backend
  is unavailable (unless \texttt{force} was specified, in which case it panics).
\item \textbf{Validate and initialize.}  \texttt{validate\_hot\_ops()} checks
  that all HOT ops (\texttt{vcpu\_run}, \texttt{mm\_region\_added},
  \texttt{mm\_region\_removed}, \texttt{context\_switch},
  \texttt{read\_clock\_ns}) are non-NULL.  Then \texttt{probe()} and
  \texttt{init()} are called on the selected backend.
\end{enumerate}

\subsection{Capability Flags}

The four boolean capability flags on the ops table replace the legacy
\texttt{using\_seccomp} global integer with structured, per-behavior queries.
This design allows code throughout UML to ask ``does the active backend use
futex-based stub dispatch?''\ rather than ``is this the seccomp backend?''---a
distinction that matters when the KVM backend might share some traits with
seccomp (e.g., futex-based communication for memory operations) but not
others (e.g., no SIGSYS filter).

\begin{center}
\begin{tabular}{lcccc}
\toprule
\textbf{Flag} & \textbf{ptrace} & \textbf{seccomp} & \textbf{KVM} \\
\midrule
\texttt{uses\_stub\_reaper}        & -- & \checkmark & -- \\
\texttt{has\_syscall\_stub\_fd\_map} & -- & \checkmark & -- \\
\texttt{stub\_syscall\_uses\_futex} & -- & \checkmark & -- \\
\texttt{stub\_child\_runs\_seccomp} & -- & \checkmark & -- \\
\bottomrule
\end{tabular}
\end{center}

(The ptrace backend has been removed; its column is included for historical
reference.  The KVM backend's flags will be populated when its ops table
matures past the stub stage.)

% ------------------------------------------------------------------
\section{Key Code Listings: Annotated Excerpts}
\label{sec:impl-code-listings}

This section presents annotated excerpts of the three most important
implementation artifacts discussed above, with commentary on the
non-obvious design choices.

\subsection{The Complete Futex Protocol in \texttt{stub\_signal\_interrupt}}
\label{sec:impl-futex-listing}

\begin{lstlisting}[language=KernelC,caption={Futex protocol in \texttt{stub\_signal\_interrupt()} (from \texttt{arch/um/kernel/skas/stub.c})},label=lst:futex-protocol-full]
/* Phase 1: Hand control to UML kernel */
restart_wait:
    d->futex = FUTEX_IN_KERN;
    do {
        res = stub_syscall3(__NR_futex,
                            (unsigned long)&d->futex,
                            FUTEX_WAKE, 1);
    } while (res == -EINTR);

    /* Phase 2: Sleep until UML kernel wakes us */
    do {
        res = stub_syscall4(__NR_futex,
                            (unsigned long)&d->futex,
                            FUTEX_WAIT, FUTEX_IN_KERN, 0);
    } while (res == -EINTR || d->futex == FUTEX_IN_KERN);

    /* Phase 3: Validate -- unexpected errors are fatal */
    if (res < 0 && res != -EAGAIN)
        stub_syscall1(__NR_exit_group, 1);

    /* Phase 4: Process any queued memory operations */
    if (d->syscall_data_len) {
        /* Receive FDs via SCM_RIGHTS over Unix socket */
        do {
            res = stub_syscall3(__NR_recvmsg, 0,
                                (unsigned long)&msghdr, 0);
        } while (res == -EINTR);
        if (res < 0 && res != -EAGAIN)
            stub_syscall1(__NR_exit_group, 1);

        /* Execute mmap/munmap batch */
        res = syscall_handler(fd_map);

        /* Close transferred FDs */
        while (num_fds)
            stub_syscall2(__NR_close, fd_map[--num_fds], 0);
    }

    /* Phase 5: Re-enter wait if batch failed or restart
     * was requested (sets signal=SIGSYS for the kernel) */
    if (res < 0 || d->restart_wait) {
        d->signal = SIGSYS;
        d->restart_wait = 0;
        goto restart_wait;
    }

    /* Phase 6: Restore architecture state and return */
    stub_seccomp_restore_state(&d->arch_data);
    /* rt_sigreturn restores guest mcontext */
\end{lstlisting}

Key design decisions in this protocol:

\begin{enumerate}[nosep]
\item \textbf{Retry on \texttt{-EINTR}.}  Both the wake and wait loops
  explicitly retry on \texttt{EINTR} because any unmasked signal can interrupt
  the futex syscall.
\item \textbf{\texttt{-EAGAIN} is benign.}  \texttt{FUTEX\_WAIT} returns
  \texttt{-EAGAIN} if the futex value changed between the comparison and the
  sleep---this is a spurious wakeup, handled by the outer check
  \texttt{d->futex == FUTEX\_IN\_KERN}.
\item \textbf{All other errors are fatal.}  The stub child operates in a
  minimal environment; unexpected futex errors indicate corruption or kernel
  bugs, and the only safe response is to terminate.
\item \textbf{SCM\_RIGHTS fd passing.}  File descriptors needed by
  \texttt{mmap} (e.g., \texttt{physmem\_fd}) cannot be directly referenced by
  the stub child (different process).  UML sends them over a Unix socket
  using the SCM\_RIGHTS ancillary message mechanism, and the stub receives
  them in the \texttt{fd\_map} array.
\item \textbf{Restart loop.}  If a batched syscall fails or the UML kernel
  sets \texttt{restart\_wait}, the stub re-enters the futex wait loop rather
  than returning to guest code.  This gives the UML kernel a chance to retry
  or take corrective action.
\end{enumerate}

\subsection{The \texttt{syscall\_handler} Batch Processor}
\label{sec:impl-syscall-handler}

\begin{lstlisting}[language=KernelC,caption={The stub's syscall batch processor},label=lst:syscall-handler]
static __always_inline int syscall_handler(int fd_map[STUB_MAX_FDS])
{
    struct stub_data *d = get_stub_data();
    int i;
    unsigned long res;
    int fd;

    for (i = 0; i < d->syscall_data_len; i++) {
        struct stub_syscall *sc = &d->syscall_data[i];

        switch (sc->syscall) {
        case STUB_SYSCALL_MMAP:
            fd = fd_map ? fd_map[sc->mem.fd] : sc->mem.fd;
            res = stub_syscall6(STUB_MMAP_NR,
                    sc->mem.addr, sc->mem.length,
                    sc->mem.prot,
                    MAP_SHARED | MAP_FIXED,
                    fd, sc->mem.offset);
            if (res != sc->mem.addr) {
                d->err = res;
                d->syscall_data_len = i;
                return -1;
            }
            break;
        case STUB_SYSCALL_MUNMAP:
            res = stub_syscall2(__NR_munmap,
                    sc->mem.addr, sc->mem.length);
            if (res) {
                d->err = res;
                d->syscall_data_len = i;
                return -1;
            }
            break;
        default:
            d->err = -95; /* EOPNOTSUPP */
            d->syscall_data_len = i;
            return -1;
        }
    }

    d->err = 0;
    d->syscall_data_len = 0;
    return 0;
}
\end{lstlisting}

On error, the handler records which syscall failed (by truncating
\texttt{syscall\_data\_len} to \texttt{i}) and the error code in
\texttt{d->err}.  The UML kernel can then determine exactly which operation
failed and what the host kernel returned.

\vspace{1em}

\noindent
Together, these data structures and code paths form the core execution engine
of User-Mode Linux.  The backend ops table provides the abstraction layer;
the stub child and futex protocol provide the trap mechanism; the TLB sync
path bridges guest and host page tables; and the signal handling infrastructure
turns POSIX signals into kernel interrupts.  The next chapter examines the
v2 redesign that adds a KVM-based backend as an alternative to this
seccomp-based architecture.
